\chapter{Diferenças em Diferenças}
\label{cap:did}
% ── índice ──────────────────────────────────────────────────────
\index{diferenças em diferenças|textbf}
\index{DiD|see{diferenças em diferenças}}
\index{did@\texttt{did()}|textbf}
\index{tendências paralelas|textbf}
\index{ATT|see{efeito médio do tratamento sobre tratados}}
\index{efeito médio do tratamento sobre tratados|textbf}

% ─────────────────────────────────────────────────────────────────────────────
\section{O modelo}

Diferenças em Diferenças (DiD) estima o efeito causal de um tratamento
comparando a variação temporal do grupo tratado com a do grupo controle:

\[
  \hat\delta_{\mathrm{DiD}}
  = \underbrace{(\bar{Y}_{1,\text{pós}} - \bar{Y}_{1,\text{pré}})}_{\text{variação tratado}}
  - \underbrace{(\bar{Y}_{0,\text{pós}} - \bar{Y}_{0,\text{pré}})}_{\text{variação controle}}
\]

A tabela $2\times 2$ resume os quatro grupos:

\begin{center}
\begin{tabular}{lcc}
\toprule
 & \textbf{Pré} & \textbf{Pós} \\
\midrule
Controle ($D=0$) & $\mu_{00}$ & $\mu_{01}$ \\
Tratado  ($D=1$) & $\mu_{10}$ & $\mu_{11}$ \\
\bottomrule
\end{tabular}
\end{center}

\[
  \hat\delta = (\mu_{11} - \mu_{10}) - (\mu_{01} - \mu_{00})
\]

% ─────────────────────────────────────────────────────────────────────────────
\section{Forma de regressão}

O estimador DiD é equivalente ao coeficiente da interação em:

\[
  y_{it} = \beta_0 + \beta_1\,\text{Post}_t + \beta_2\,\text{Treated}_i
           + \delta\,(\text{Treated}_i \times \text{Post}_t) + \varepsilon_{it}
\]

onde $\hat\delta = \hat\beta_3$ é o efeito médio do tratamento sobre os
tratados (ATT — \emph{Average Treatment effect on the Treated}).

% ─────────────────────────────────────────────────────────────────────────────
\section{Hipótese de tendências paralelas}

A identificação do ATT repousa sobre:

\[
  E\!\left[Y_{it}^{(0)} - Y_{is}^{(0)} \mid D_i = 1\right]
  = E\!\left[Y_{it}^{(0)} - Y_{is}^{(0)} \mid D_i = 0\right]
  \qquad \forall\; t \neq s
\]

Na ausência do tratamento, tratados e controles teriam a \emph{mesma}
tendência temporal. A suposição é não testável no período pós-tratamento,
mas pode ser verificada nos períodos pré-tratamento quando há mais de um
período disponível.

\begin{aviso}
  A hipótese de tendências paralelas é mais forte do que parece. Grupos
  com tendências de crescimento estruturalmente diferentes (ex.\ regiões
  em expansão vs.\ regiões estagnadas) violam a hipótese mesmo que os
  níveis sejam similares no pré-tratamento.
\end{aviso}

% ─────────────────────────────────────────────────────────────────────────────
\section{Verificação por DGP}

\subsection*{Recuperação exata (sem ruído)}

O DGP usa 500 indivíduos × 2 períodos ($n = 1000$). Os primeiros 250
são tratados, os demais são controle. O efeito verdadeiro é $\delta = 3{,}0$:

\[
  y_{it} = 1{,}0 + 0{,}5\cdot\text{Post}_t + 3{,}0\cdot(D_i\times\text{Post}_t)
\]

As quatro médias de célula são $\mu_{00}=1{,}0$, $\mu_{01}=1{,}5$,
$\mu_{10}=1{,}0$ e $\mu_{11}=4{,}5$, de modo que
$\hat\delta = (4{,}5 - 1{,}0) - (1{,}5 - 1{,}0) = 3{,}0$ exatamente.

\begin{lstlisting}[caption={DGP DiD sem ruído — recuperação exata}]
seed(42)
input df
id
1
end
for i in 1..=10 { df = append(df, df) }
df |> filter(_n <= 1000)
generate df id      = (_n - 1) % 500 + 1
generate df post    = _n > 500
generate df treated = id <= 250
generate df att     = treated * post
generate df y       = 1.0 + 0.5*post + 3.0*att
did(y ~ treated + post, df)
\end{lstlisting}

\begin{lstlisting}[language={}, basicstyle=\ttfamily\small,
  frame=none, numbers=none, backgroundcolor=\color{codbg},
  xleftmargin=0pt, xrightmargin=0pt, breaklines=false]
================= Difference-in-Differences (2x2 Canonical) ==================
ATT (Effect):                 3.0000 || R-squared:                    1.0000
Std. Error:                   0.0000 || P-value:                      0.0000
t-statistic:                     inf || Observations:                   1000

-------------------------------- Group Means ---------------------------------
Control Group (Pre):      1.0000 | Control Group (Post):     1.5000
Treated Group (Pre):      1.0000 | Treated Group (Post):     4.5000
Parallel Trend Diff:      0.0000 (If > 0, Control grew more than Treated Pre-trend)
==============================================================================
\end{lstlisting}

\subsection*{Com ruído idiossincrático}

Adicionando $\varepsilon \sim \mathcal{N}(0,1)$, o DiD permanece consistente:

\begin{lstlisting}[caption={DGP DiD com ruído}]
generate df e = rnormal()
generate df y = 1.0 + 0.5*post + 3.0*att + e
did(y ~ treated + post, df)
\end{lstlisting}

\begin{lstlisting}[language={}, basicstyle=\ttfamily\small,
  frame=none, numbers=none, backgroundcolor=\color{codbg},
  xleftmargin=0pt, xrightmargin=0pt, breaklines=false]
================= Difference-in-Differences (2x2 Canonical) ==================
ATT (Effect):                 3.0114 || R-squared:                    0.9645
Std. Error:                   0.0358 || P-value:                      0.0000
t-statistic:                 84.0942 || Observations:                   1000

-------------------------------- Group Means ---------------------------------
Control Group (Pre):      1.4842 | Control Group (Post):     2.0099
Treated Group (Pre):      1.4880 | Treated Group (Post):     5.0251
Parallel Trend Diff:      0.0000 (If > 0, Control grew more than Treated Pre-trend)
==============================================================================
\end{lstlisting}

O ATT de $3{,}0114$ está próximo do verdadeiro $3{,}0$; o desvio é variância
amostral. O \emph{Parallel Trend Diff} próximo de zero confirma que a
hipótese de tendências paralelas é satisfeita por construção.

% ─────────────────────────────────────────────────────────────────────────────
\section{Sintaxe}

\begin{lstlisting}
did(outcome ~ treated + post, df)
did(outcome ~ treated + post, df, cov=robust)
\end{lstlisting}

A fórmula requer exatamente duas variáveis no lado direito: a indicadora de
grupo (\lstinline{treated}) e a indicadora de período (\lstinline{post}).
O \hayashi{} calcula a interação internamente.

\begin{lstlisting}[caption={Fluxo típico de DiD}]
load "dados.csv" as df

// cria indicadoras
generate df post    = ano >= 2020
generate df treated = estado == "RS"

let m = did(salario ~ treated + post, df, cov=robust)
display m
\end{lstlisting}

\section{Capturando o resultado}

\begin{lstlisting}
let m = did(y ~ treated + post, df)
display m
\end{lstlisting}

\begin{nota}
  Para estudos de evento com múltiplos períodos pré e pós-tratamento,
  o DiD canônico $2\times 2$ pode ser insuficiente. Nesses casos,
  estime o modelo de regressão com interações período a período e
  verifique a ausência de efeitos pré-tratamento como teste de
  tendências paralelas.
\end{nota}

\section{Event study (DiD dinâmico)}

\index{eventstudy@\texttt{eventstudy()}}
\index{event study}

O comando \texttt{eventstudy} estima um DiD dinâmico com leads e lags,
construindo dummies de tempo de evento para cada período relativo ao
evento (excluindo o período de referência):

\begin{lstlisting}
eventstudy(y ~ event_time + controles, df, ref=-1, min=-5, max=5, cov=HC1)
\end{lstlisting}

\begin{itemize}
  \item \texttt{event\_time}: variável com o tempo relativo ao evento
    ($\ldots, -2, -1, 0, 1, 2, \ldots$)
  \item \texttt{ref}: período de referência omitido (padrão $-1$)
  \item \texttt{min}, \texttt{max}: limites inferior e superior dos
    tempos de evento incluídos (padrão $-5$ e $5$)
  \item Controles opcionais são incluídos após \texttt{event\_time}
\end{itemize}

O resultado reporta coeficientes, erros-padrão, estatísticas $t$ e
$p$-values para cada tempo de evento. Os coeficientes dos leads
($t < 0$) testam a hipótese de tendências paralelas: se
estatisticamente significativos, indicam violação do pressuposto.

% ─────────────────────────────────────────────────────────────────────────────
\section{Local Projections Difference-in-Differences}
\label{sec:lpdid}

\index{lpdid@\texttt{lpdid()}}
\index{Local Projections Difference-in-Differences}

O comando \texttt{lpdid} implementa o estimador de
Dube, Girardi, Jordà e Taylor (2025), que combina projeções locais com
Diferenças-em-Diferenças para designs de adoção escalonada. Em vez
de dummies de tempo de evento, o LP-DiD estima o ATT em cada horizonte
$h$ por uma regressão de diferenças longas:
\[
  \Delta_h y_{it} = \tau_h \, \Delta_h D_{it} + \text{controles} + \varepsilon_{it},
\]
onde $\Delta_h y_{it} = y_{i,t+h} - y_{i,t-\ell}$ é a diferença
longa do outcome, $\ell$ é o período base e $D_{it}$ indica a
unidade tratada. A amostra em cada horizonte é restrita a unidades
``limpas'' (clean controls): nunca tratadas ou ainda não tratadas em
$t+h$, conforme a opção \texttt{clean\_control}.

\begin{lstlisting}[caption={LP-DiD com adoção escalonada absorvente}]
let m = lpdid(y ~ unit + time + first_treat, df,
              max_pre=5, max_post=8, base_period=-1,
              clean_control="not_yet_treated")
print(m.event_coefficients)
\end{lstlisting}

Os principais argumentos são:

\begin{itemize}
  \item \texttt{outcome \textasciitilde\ unit + time + first\_treat}:
    fórmula com outcome, identificador da unidade, tempo e primeiro
    período de tratamento (0 indica nunca tratado).
  \item \texttt{max\_pre}, \texttt{max\_post}: número de horizontes
    pré e pós-tratamento a estimar.
  \item \texttt{base\_period}: período base para a diferença longa.
    Pode ser um inteiro negativo (ex. \texttt{-1}), uma lista de
    inteiros negativos ou \texttt{"all\_pre"}.
  \item \texttt{clean\_control}: estratégia de controle limpo
    (\texttt{not\_yet\_treated}, \texttt{never\_treated},
    \texttt{first\_entry}, \texttt{stabilized}).
  \item \texttt{estimand}: estimando alvo \texttt{vw} (padrão,
    variance-weighted), \texttt{rw} (equally-weighted via
    Frisch-Waugh-Lovell) ou \texttt{ra} (regression-adjusted
    imputation).
  \item \texttt{nonabsorbing}: permite que o tratamento ligue e
    desligue ao longo do tempo. Nesse caso use a variável de
    tratamento binário \texttt{treatment=} em vez de
    \texttt{first\_treat}.
  \item \texttt{fixed\_composition}: mantém o mesmo conjunto de
    unidades em todos os horizontes, exigindo controles não tratados
    em $t+H$.
\end{itemize}

O objeto de resultado possui três campos principais:
\texttt{event\_coefficients} (tabela horizonte a horizonte),
\texttt{scalar\_coefficients} (ATT médio e ATT pooled) e
\texttt{fit} (metadados do ajuste).
